% ---- pg-418 （书页 406） ----

\begin{mvfigure}
\mvinc{pg-418}{0.314,0.124,0.685,0.337}
\phantomsection\label{fig:9.5}\mvcap{图 9.5\quad 鸡的骨骼数据前两个因子的因子得分}
\end{mvfigure}

画出用主成分和极大似然估计因子载荷而得到因子得分配对的图形也是有意义的. 对于鸡的骨骼数据, 因子得分配对的图示于\hyperref[fig:9.6]{图 9.6}. 如果在某个因子上诸载荷一致, 则此得分配对将紧紧聚集在通过原点的 $45^\circ$ 直线附近. 不一致的那些载荷, 将产生偏离这一模式的因子得分. 如果后一情况出现, 它通常与最后面的因子有关, 并且可能表明因子数太大了, 即最后的那些因子没有多大意义. 这似乎正是鸡的骨骼数据中第三因子出现的情况, 如\hyperref[fig:9.6]{图 9.6} 的(c)图所示.

\begin{mvfigure}
\mvinc{pg-418}{0.218,0.532,0.757,0.795}
\mvcap{(a) 第一因子}
\phantomsection\label{fig:9.6}\mvcap{图 9.6\quad 鸡的骨骼数据的因子得分配对\par（用主成分法和极大似然法估计载荷）}
\end{mvfigure}
